\section{One original native frame transition controls every degree}

\begin{definition}\label{def:native-tc-compatible-product-residue-frames}
\lean{EnvelopeProperty.NativeTestConfiguration.TestConfiguration.productReferenceResidueFrame,EnvelopeProperty.NativeTestConfiguration.TestConfiguration.productModelResidueFrame,EnvelopeProperty.NativeTestConfiguration.TestConfiguration.productModelPowerResidueFrame,EnvelopeProperty.NativeTestConfiguration.TestConfiguration.productFrameTransitionUnit}
\leanok
\uses{def:scheme-support-line-frame,def:native-tc-product-original-line-iso,def:fc-induced-native-residue-frame,def:fc-original-frame-transition-unit}
Fix an actual native point $y$ of the original product, an actual
original open $U\subseteq X$ containing its projected support,
an original frame $e_U$ of $L$ on $U$, and an actual residue frame
$e_M$ of the original total line at the image of $y$. The
reference residue frame $e_{\rm ref}$ is the actual support
transport of the local product pullback of $e_U$. The model
residue frame $e_{\rm mod}$ is the native induced frame through
the original whole product line isomorphism. In degree $n$ use
the original power frame of $e_M$ and the literal product power
isomorphism to obtain $e_{{\rm mod},n}$. Define the single native
unit $u$ to be the transition from $e_{\rm ref}$ to $e_{\rm mod}$.
\end{definition}
\begin{proof}
The actual support point lies in the inverse-image open, so its
native residue morphism factors through that open. The local
product pullback of the actual original frame is a whole frame
there; support transport therefore defines $e_{\rm ref}$. The
original action and fibre-one data define the actual product
line isomorphism, and native residue transport of $e_M$ through
it defines $e_{\rm mod}$. Apply the same construction with the
literal power isomorphism and the genuine pulled power frame of
$e_M$ for each $n$. Two whole frames of the same native residue
line differ by the actual unit obtained from their composite
automorphism of the native structure sheaf. This is $u$, fixed
before all degrees and all section evaluations.
\end{proof}

\begin{theorem}\label{thm:native-tc-compatible-product-power-frames}
\lean{EnvelopeProperty.NativeTestConfiguration.TestConfiguration.productTotalLinePowerIso_inducedResidueFrame,EnvelopeProperty.NativeTestConfiguration.TestConfiguration.productReferenceResidueFrame_power}
\leanok
\uses{def:native-tc-compatible-product-residue-frames,thm:native-tc-literal-product-power-coherence,cor:actual-induced-normalized-power-frame,thm:actual-support-local-pullback-power-frame}
The reference frame in degree $n$, constructed from the same
original $U,e_U$ by actual powers, and the actual model power
frame each equal one common native residue pullback of the
canonical product-projection power comparison, followed by the
genuine pulled power frame of their corresponding original
degree-one residue frame.
\end{theorem}
\begin{proof}
The literal product power coherence identifies the original
power isomorphism with the normalized canonical power of the
original product line isomorphism. The actual induced normalized
power theorem therefore identifies the model power frame with
the stated common comparison followed by the pulled power of
$e_{\rm mod}$. The support local pullback power theorem, applied
to the actual product projection, identifies the reference power
frame with exactly that same comparison followed by the pulled
power of $e_{\rm ref}$. Both comparisons use the unchanged native
residue morphism and the unchanged original product projection.
\end{proof}

\begin{theorem}\label{thm:native-tc-frame-transition-unit-powers}
\lean{EnvelopeProperty.NativeTestConfiguration.TestConfiguration.productPowerFrame_transitionUnit,EnvelopeProperty.NativeTestConfiguration.TestConfiguration.productFrameTransitionUnit_value_pos,EnvelopeProperty.NativeTestConfiguration.TestConfiguration.productPowerFrame_transitionUnit_value}
\leanok
\uses{thm:native-tc-compatible-product-power-frames,thm:fc-actual-power-frame-transition-units,lem:fc-original-frame-transition-positive-value}
The actual native transition unit from the original reference
power frame to the original model power frame is $u^n$ for all
$n\geq0$. For every valuation $\nu$ of the actual product residue
field, $c=\nu(u)>0$, and the value of this degree-$n$ transition
unit is $c^n$. The unit and positive factor are fixed before
every degree, original section and original finite family.
\end{theorem}
\begin{proof}
Replace the two degree-$n$ frames using the preceding theorem.
Transition units are unchanged when both frames are preceded
by the same whole isomorphism, so cancel their common native
product-projection power comparison. Expand their two pulled
power frames and cancel the common native residue power pullback
comparison as well. The remaining frames are the genuine entire
power frames of $e_{\rm ref}$ and $e_{\rm mod}$. Their actual unit
transition is the power of the original transition unit, hence
$u^n$. A valuation is positive on a unit because a unit and its
inverse multiply to one, and it carries powers to powers.
Applying this to the actual unit equality proves the value
identity, including degree zero. No global section in degree one
is used to define the factor or establish its positivity.
\end{proof}
